% !TEX root = /home/eb/.vscode/extensions/qftrules.programme-colle/tmp/Exercice.tex
%%%%%%%%%%%%%%%%%%%%%%%%%%%%%%%%%%%%%%%%%%%%%%%%%%%%%%%%%%%%%%%

%%%%%%%%%%%%%%%%%%%%%%%%%%%%%%%%%%%%%%%%%%%%%%%%%%%%%%%%%%%%%%%%%%%%%%%%%%%%%%%%%%%%%%%
\begin{exo}[PC][2][TD]{Couche anti-reflet}
  \Notions{\item Onde électromagnétique \task Indice optique}
À la surface d'un verre d'indice $n$ on dépose une couche anti-reflet d'épaisseur $e$ et d'indice $N$ de telle sorte que l'espace est divisé en 3 domaines: l'air ($x<0$) d'indice 1, la couche anti-reflet $(0<x<e$) d'indice $N$ et le verre $(x>e)$ d'indice $n$. Une OPPH de pulsation $\omega$ polarisée rectilignement sur $\vv{e_z}$ arrive dans l'air sous incidence normale et on cherche les conditions pour qu'il n'y ait pas d'onde réfléchie dans l'air.
 \begin{quest}
\item Donner en notation complexe la forme attendue du champ électrique dans chaque région.
\sol{
On suppose qu'il n'y a pas d'onde réfléchie dans l'air.\\
Dans l'air : champ incident $\udl{\E_0}=\udl{E_{i}} e^{i(\omega t - k_0 x)}\ez$.\\
Dans la couche : superposition des champs transmis et réfléchis : $\udl{\E_1}=\udl{E_{t1}} e^{i(\omega t - k_1 x)}\ez + \udl{E_{r1}} e^{i(\omega t + k_1 x)}\ez$.\\
Dans le verre : onde transmise uniquement : $\udl{\E_2}=\udl{E_{t}} e^{i(\omega t - k_2 (x-e))}\ez$.\\
On a $k_0 = \frac{\omega}{c}$, $k_1 = N \frac{\omega}{c}$ et $k_2 = n \frac{\omega}{c}$.
}
\item En déduire la forme du champ magnétique dans les trois milieux.
\sol{
Relation de structure pour une OPPH : $\Vec{B} = \frac{\Vec{k}}{\omega} \pv \E = \frac{n}{c} \Vec{u} \pv \E$ avec $\Vec{u} = \pm \ex$ suivant la direction de propagation. Attention à bien appliquer cette formule sur chaque OPPH séparément (et pas directement sur $\udl{\E_1}$).
}
\item Établir 4 relations entre les amplitudes du champ électrique dans les trois régions, puis montrer que
\begin{equation*}
\frac{(N+1)(N-n)}{(N-1)(N+n)}=\e^{i2Nk_{0}e},
\end{equation*}
où $k_{0}$ est le module du vecteur d'onde dans le vide.
\sol{
On utilise la continuité des champs électrique et magnétique à chaque interface. On repère l'interface air-couche en $x=0$. On obtient alors
$$
\begin{cases}
\udl{E_{i}} = \udl{E_{t1}} + \udl{E_{r1}} \\
\udl{E_{i}} = N(\udl{E_{t1}} - \udl{E_{r1}}) \\
\end{cases}\qquad
\begin{cases}
\udl{E_{t1}} e^{-i N k_0 e)} + \udl{E_{r1}} e^{+i N k_0 e)} = \udl{E_{t}}\\
N(\udl{E_{t1}} e^{-i N k_0 e)} - \udl{E_{r1}} e^{+i N k_0 e)}) = n\udl{E_{t}} \\
\end{cases}
$$
Le premier système donne $\udl{E_{t1}}=\frac{N+1}{2N}\udl{E_{i}}$ et $\udl{E_{r1}}=\frac{N-1}{2N}\udl{E_{i}}$. On en déduit
$$
\begin{cases}
\udl{E_{i}}(\frac{N+1}{2N} e^{-i N k_0 e)} + \frac{N-1}{2N} e^{+i N k_0 e)}) = \udl{E_{t}}\\
\udl{E_{i}}(\frac{N+1}{2} e^{-i N k_0 e)} - \frac{N-1}{2} e^{+i N k_0 e)}) = n\udl{E_{t}} \\
\end{cases}
$$
Or la conservation de l'énergie impose $\udl{E_{i}} = \udl{E_{t}}$, et donc on trouve la relation attendue.
}
\item En déduire les valeurs qu'il faut choisir pour $e$ et $N$. Comment expliquer qu'il subsiste des reflets violet et rouge sur un objectif ainsi traité lorsqu'on l'observe à la lumière du Soleil ?
\sol{
En l'absence d'absorption, $\frac{(N+1)(N-n)}{(N-1)(N+n)}$ est réel et donc $\e^{i2Nk_{0}e}=\pm 1$, soit $2Nk_0 e = p\pi$ avec $p$ entier, et donc
\eqbox{eN = p \frac{\lambda_0}{4}.}
Cette condition dépend de la longueur d'onde $\lambda_0$ (et dans une moindre mesure, l'indice optique $N$ aussi). En fonction du signe de $\e^{i2Nk_{0}e}$, on trouve alors deux valeurs possibles pour $N$ (l'une supérieure à $n$ et l'autre inférieure) puis pour $e$, données par l'équation ci-dessus.
}
\end{quest}
\end{exo}
%%%%%%%%%%%%%%%%%%%%%%%%%%%%%%%%%%%%%%%%%%%%%%%%%%%%%%%%%%%%%%%%%%%%%%%%%%%%%%%%%%%%%%%


%%%%%%%%%%%%%%%%%%%%%%%%%%%%%%%%%%%
\begin{exo}[PC][3][TD]{Coefficients de Fresnel}
  \Notions{\item Onde électromagnétique \task Indice optique}
Une OPPH électromagnétique de pulsation $\omega$ arrive dans l'air sous un angle d'incidence $\theta_i$ sur une surface de verre d'indice optique réel $n$. On note $\theta_r$ l'angle de réfraction. La normale au dioptre est portée par l'axe $(Oy)$ : le plan d'incidence est alors repéré par les axes $(Ox)$ et $(Oy)$, tandis que le plan du dioptre est repéré par les axes $(Ox)$ et $(Oz)$, comme sur la \Fig{Fresnel_coeff}. Le champ électrique est supposé de polarisation rectiligne, et d'expression complexe 
\begin{equation}
\udl{\E} = \Vec{E_0} \,e^{i(\omega t - \vec{k}\cdot\vec{r})}.
\end{equation}
On cherche à calculer les coefficients de réflexion et de transmission en amplitude, appelés coefficients de Fresnel. On suppose pour commencer l'onde est polarisée dans le plan du dioptre, soit 
$\Vec{E_0} = E_0 \, \ez$.
\begin{fig}[Fresnel_coeff]
  \centering
  \includegraphics[width=0.4\columnwidth]{fresnel.pdf}
  %\capt{}
\end{fig}
 \begin{quest}
\item Donner en notation complexe la forme attendue du champ électrique dans chaque région, air et verre. On introduira le vecteur d'onde $k_0$ dans l'air, assimilé au vide.
\sol{
  La réflexion et la réfraction conserve la polarisation.
  \begin{enumerate}
    \item Dans l'air.
    $$
    \begin{cases}
\text{Champ incident : } \udl{\E}_i={E}_{0,i} e^{i(\omega t - k_0 (x\sin \theta_i + y \cos \theta_i))}\ez\\
\text{Champ réfléchi : } \udl{\E}_r={E}_{0,r} e^{i(\omega t - k_0 (x\sin \theta_i - y \cos \theta_i))}\ez
\end{cases}
$$
    \item Dans le verre.
    $$
    \text{Champ transmis : } \udl{\E}_t={E}_{0,t} e^{i(\omega t - n k_0 (x\sin \theta_r + y \cos \theta_r))}\ez
$$
  \end{enumerate}
% Dans l'air : champ incident $\udl{\E_0}=\udl{E_{i}} e^{i(\omega t - k_0 (x\sin \theta_i + y \cos \theta_i))}\ez$ et champ réfléchi $\udl{\E_r}=\udl{E_{r}} e^{i(\omega t - k_0 (x\sin \theta_i - y \cos \theta_i))}\ez$.\\
% Dans le verre : champ transmis $\udl{\E_t}=\udl{E_{t}} e^{i(\omega t - n k_0 (x\sin \theta_r + z \cos \theta_r))}\ez$.\\
% On a $k_0 = \frac{\omega}{c}$ et $k_2 = n \frac{\omega}{c}$, et $\theta_r$ est l'angle de réfraction donné par la loi de Snell-Descartes : $n \sin \theta_r = \sin \theta_i$.
}
\item En déduire la forme du champ magnétique dans les chaque région.
\sol{La relation de structure pour une OPPH est 
\eqbox{
\Vec{B} =  \frac{n}{c} \Vec{u} \pv \E
}
avec $\Vec{u}$ le unitaire dans la direction de propagation. On en déduit alors
\begin{enumerate}
  \item Dans l'air.
  $$
  \begin{cases}
\text{Champ incident : } \udl{\B}_i=\frac{1}{c} \left(\sin \theta_i \ex - \cos \theta_i \ey \right) \pv \udl{\E}_i =  \frac{1}{c} \left(\sin \theta_i \ex - \cos \theta_i \ey \right) E_{0,i} e^{i(\omega t - k_0 (x\sin \theta_i + y \cos \theta_i))}\ey\\
\text{Champ réfléchi : } \udl{\B}_r=\frac{1}{c} \left(\sin \theta_i \ex + \cos \theta_i \ey\right) \pv \udl{\E}_r = \frac{1}{c} \left(\sin \theta_i \ex + \cos \theta_i \ey\right) E_{0,r} e^{i(\omega t - k_0 (x\sin \theta_i - y \cos \theta_i))}\ey
\end{cases}
$$
  \item Dans le verre.
  $$
  \text{Champ transmis : } \udl{\B}_t=\frac{n}{c} \left(\sin \theta_r \
ex - \cos \theta_r \ey\right) \pv \udl{\E}_t = \frac{n}{c} \left(\sin \theta_r \ex - \cos \theta_r \ey\right) E_{0,t} e^{i(\omega t - n k_0 (x\sin \theta_r + y \cos \theta_r))}\ey
$$
\end{enumerate}
} 
\item Définir les coefficients de réflexion et de transmission en amplitude, $r$ et $t$, puis établir 4 relations entre les grandeurs $r$, $t$, $\theta_i$ et $\theta_r$. Montrer la loi de Snell-Descartes à partir de ces relations.
  \sol{
On utilise la continuité des champs électrique et magnétique à l'interface $y=0$. On obtient alors
$$\begin{cases}
(E_{0,i} + E_{0,r})e^{i(\omega t - k_0 x \sin \theta_i)} = E_{0,t} e^{i(\omega t - n k_0 x \sin \theta_r)}\\
\frac{1}{c} (E_{0,i} (\sin \theta_i \ex - \cos \theta_i \ey) + E_{0,r} (\sin \theta_i \ex + \cos \theta_i \ey)) e^{i(\omega t - k_0 x \sin \theta_i)} = \frac{n}{c} E_{0,t} (\sin \theta_r \ex - \cos \theta_r \ey) e^{i(\omega t - n k_0 x \sin \theta_r)}
\end{cases}$$ 
On projette ces relations dans le plan d'incidence $(Ox,Oy)$ et on trouve alors
$$\begin{cases}
E_{0,i} + E_{0,r} = E_{0,t}\\
\sin \theta_i = n \sin \theta_r\\
E_{0,i} \sin \theta_i + E_{0,r} \sin \theta_i = n E_{0,t} \sin \theta_r\\
- E_{0,i} \cos \theta_i + E_{0,r} \cos \theta_i = - n E_{0,t} \cos \theta_r
\end{cases}$$
En faisant apparaitre les coefficients de réflexion et de transmission en amplitude, $r = \frac{E_{0,r}}{E_{0,i}}$ et $t = \frac{E_{0,t}}{E_{0,i}}$, on trouve alors les relations suivantes :
\eqbox{
\begin{cases}
1 + r = t\\
\sin \theta_i = n \sin \theta_r\\
(1 + r) \sin \theta_i = n t \sin \theta_r\\
- (1 - r) \cos \theta_i = - n t \cos \theta_r
\end{cases}
}
  }
\item En déduire les coefficients de réflexion et de transmission en amplitude, $r$ et $t$, en fonction de $\theta_i$, $\theta_r$, puis en fonction de $\theta_i$ uniquement. Tracer l'allure des graphes $r(\theta_i)$ et $t(\theta_i)$ pour $n=\num{1.41}$.  Ces coefficients peuvent-ils s'annuler ? Interpréter.
  \sol{On trouve
\eqbox{
r = \frac{\cos \theta_i - n \cos \theta_r}{\cos \theta_i + n \cos \theta_r} \qquad\text{et}\qquad t = \frac{2 \cos \theta_i}{\cos \theta_i + n \cos \theta_r}
}
  }
\end{quest}
On considère maintenant que l'onde est polarisée perpendiculairement au plan du dioptre, soit $\Vec{E_0} = E_0 \, \ex$. On admet que les coefficients de réflexion et de transmission sont alors donnés par les formules suivantes :
\begin{equation}
r = \frac{n \cos \theta_i - \cos \theta_r}{n \cos \theta_i + \cos \theta_r} \qquad\text{et}\qquad t = \frac{2 n \cos \theta_i}{n \cos \theta_i + \cos \theta_r}
\end{equation}
\begin{quest}
  \item	Montrer que l'onde réfléchie s'annule pour une incidence particulière, appelée angle de Brewster $\theta_B$. Exprimer $\theta_B$ en fonction de $n$, puis calculer sa valeur pour $n=\num{1.41}$.
\end{quest}
\end{exo}
%%%%%%%%%%%%%%%%%%%%%%%%%%%%%%%%%%%